% GenAIDeclarationPaper — technical report (skeleton)
% Open source at https://github.com/ceurws/GenAIDeclarationPaper
%
% Target: technical report on Zenodo. Not peer-reviewed.
% Experiment design: https://cr.bitplan.com/index.php/CEUR-WS_Declaration_on_Generative_AI_Check
%
% CEUR-WS 1-column style (ceurart). Do NOT use the twocolumn / hf options.
% The report follows the CEUR-WS Policy on AI-Assisting Tools although it is
% not published via CEUR-WS; see AGENTS.md.
%
% Sections are intentionally EMPTY until the authors write them.
% The report must be self-contained: every term, class, research question
% and method step is defined in the text; links only point to further material.
%
\documentclass{ceurart}
% typographical quoting
\newcommand{\tquote}[1]{``#1''}
% footnote links
\newcommand{\footnotelink}[2]{\href{#1}{#2}\footnote{#1}}
\sloppy

\begin{document}

%% Rights management information (CC-BY is the CEUR-WS default)
\copyrightyear{2026}
\copyrightclause{Copyright for this paper by its authors.
  Use permitted under Creative Commons License Attribution 4.0
  International (CC BY 4.0).}

%% Conference information
\conference{Technical report, 2026}

%% Title
% TODO(authors): write the title.
\title{Title to be written by the authors}
\tnotemark[1]
\tnotetext[1]{Technical report. This version has not been peer-reviewed.}

%% Authors
% Confirmed authors only. ORCIDs taken from the media wiki scholar pages.
% Further co-authors are NOT listed until participation is confirmed.
% AI agents are never authors (CEUR-WS HOWTOSUBMIT FAQ 15).
\author[1,2]{Wolfgang Fahl}[%
orcid=0000-0002-0821-6995,
email=wf@bitplan.com,
]
\cormark[1]
\address[1]{RWTH Aachen University, Computer Science i5, Aachen, Germany}
\address[2]{BITPlan GmbH, Willich, Germany}
\cortext[1]{Corresponding author.}

%% Abstract (author-provided text).
\begin{abstract}
% TODO(authors)
\end{abstract}

%% Keywords
\begin{keywords}
% TODO(authors)
\end{keywords}

\maketitle

\section{Introduction}\label{sec:introduction}
% TODO(authors)

\section{Background}\label{sec:background}
% TODO(authors): CEUR-WS Policy on AI-Assisting Tools, GenAI Usage Taxonomy,
% ceurtools and row 0 of the editor agreement; define every term used later.

\section{Prior Art}\label{sec:priorart}
% TODO(authors)

\section{Research Questions and Hypotheses}\label{sec:rq}
% TODO(authors): RQ1 to RQ4 and the classes A to D in the authors' own words.

\section{Data}\label{sec:data}
% TODO(authors): population, sample, seed, strata, weights.
% No CEUR-WS PDFs are redistributed; see data/README.md.

\section{Method}\label{sec:method}
% TODO(authors): heading search, ceurtools baseline, local language model,
% annotation and agreement, comparison; the agentic work and the provenance
% notation of the repository are described here, in the text.

\section{Results}\label{sec:results}
% TODO(authors)

\section{Self-Check}\label{sec:selfcheck}
% TODO(authors): verdict of check-pdf-errors on this report's own PDF,
% version pinned by date and checksum.

\section{Threats to Validity}\label{sec:threats}
% TODO(authors)

\section{Conclusion}\label{sec:conclusion}
% TODO(authors)

%% Declaration on Generative AI (required by CEUR-WS, see
%% https://ceur-ws.org/GenAI/Policy.html and https://ceur-ws.org/GenAI/Taxonomy.html).
%% Finalized by the authors from PROVENANCE.md before release.
%% Must name every tool and map each contribution to a taxonomy category.
%% A language model used as a research instrument (RQ4) is described in
%% Section Method and named here as such, not as writing assistance.
\section*{Declaration on Generative AI}
% TODO(authors)

%% Bibliography
\nocite{ceurws2024genaipolicy,ceurws2024genaitaxonomy}
\bibliography{references}

\end{document}
